% MG05b. Después de la inversión en la fibra de frontera y antes de
% subsec:alpha-completa; requiere tab:alpha-ventana y eq:alpha-telescopia.
\subsubsection{Controles de orientación, frontera y posición del registro}
\label{mg05:controles}

La clausura entera permite distinguir tres dependencias operativas:
sentido de propagación del acarreo, condición de borde y posición del
registro dodecafásico. Su contraste utiliza la misma ventana de doce
bloques de la tabla~\ref{tab:alpha-ventana}. Se mantienen las palabras
regionales \(P,E,\Phi\) y el registro \(K\) ya generado, salvo la
modificación identificada en cada caso. La referencia común es
\(\mathcal C_{12}\) con \(c_{13}=0\), cuya salida es
\eqref{eq:alpha-ventana-cero}.

La producción prospectiva del registro firmado y su inversión
\(K^{(r)}=H_4U^{(r)}/4\) preceden a estos controles. Se conserva la
normalización del lema~\ref{lem:k-hadamard}, en particular
\(H_4^2=4I_4\) y \(\det H_4=-16\). La inversión determina el
registro una vez producido \(U\); su selección operacional permanece
en la construcción que antecede a la clausura. Las variaciones de
\(K\) que siguen tienen la función de examinar la respuesta de un
lector ya definido.

\begin{proposition}[Separación de las dependencias de la clausura]
\label{mg05:dependencias}
Para los datos de la tabla~\ref{tab:alpha-ventana}, las siguientes
cuatro modificaciones tienen efectos distintos:
\begin{enumerate}
\item La propagación de cocientes de izquierda a derecha, iniciada con
cociente nulo, cambia el tercer bloque de \(352\) a \(353\).
\item La condición \(c_{13}=1\), con la orientación original,
cambia el último bloque de \(663\) a \(664\) y conserva los once
bloques anteriores.
\item La rotación \(K'_j=K_{j+1}\), con \(K_{13}=K_1\), cambia
la palabra de salida aun manteniendo \(c_{13}=0\).
\item La perturbación \(K'_1=K_1+1=235\), con las otras once
coordenadas fijas, cambia el primer bloque de \(007\) a \(006\)
y conserva los restantes.
\end{enumerate}
\end{proposition}
\begin{proof}
Para el primer caso, definimos el cociente inicial \(b_0=0\) y
aplicamos, en el orden \(j=1,\ldots,12\),
\[
 \widetilde s_j=P_j+E_j-\Phi_j-K_j+b_{j-1},\qquad
 \widetilde A_j=[\widetilde s_j]_{1000},\qquad
 b_j=\frac{\widetilde s_j-\widetilde A_j}{1000}.
\]
Esta orientación produce, en dos mitades consecutivas,
\begin{align*}
 \widetilde A_{1:6}&=(007,297,353,570,284,800),\\
 \widetilde A_{7:12}&=(995,285,106,471,379,663).
\end{align*}
La diferencia ya aparece en \(j=3\), donde el cociente que necesita
la clausura original procede de la cuarta posición.

Para el segundo caso, la última suma pasa de \(663\) a \(664\).
Ambas tienen cociente euclídeo cero; por ello \(c_{12}=0\) en los
dos casos y la recurrencia coincide desde \(j=11\) hasta \(j=1\).
La palabra queda
\begin{align*}
 A^{(1)}_{1:6}&=(007,297,352,569,283,800),\\
 A^{(1)}_{7:12}&=(997,285,105,472,380,664).
\end{align*}
La identidad \eqref{eq:alpha-telescopia} registra exactamente la
variación unitaria del extremo derecho.

La rotación del tercer caso, evaluada con el mismo borde nulo y la
recurrencia original, da
\begin{align*}
 A^{\mathrm{rot}}_{1:6}&=(698,699,763,639,119,004),\\
 A^{\mathrm{rot}}_{7:12}&=(559,429,226,119,926,030).
\end{align*}
En este caso el acarreo izquierdo es \(c_1=-1\); el extremo derecho
permanece nulo. La conservación de ambos extremos en
\eqref{eq:alpha-telescopia} mantiene la igualdad entera al cambiar la
posición de los doce coeficientes.

Finalmente, la perturbación de \(K_1\) deja intacta toda la
recurrencia desde \(j=12\) hasta \(j=2\). En \(j=1\), la suma
disminuye de siete a seis, con cociente cero en ambos casos. Resulta
\begin{align*}
 A^{\mathrm{pert}}_{1:6}&=(006,297,352,569,283,800),\\
 A^{\mathrm{pert}}_{7:12}&=(997,285,105,472,380,663).
\end{align*}
\end{proof}

Los cuatro controles separan el sentido de transporte, la fibra de
frontera, la posición de los coeficientes y la respuesta local. Su
interpretación conserva el dominio finito \(N=12\). El transporte
de la cola y la normalización de la sección completa se desarrollan
en la continuación, con las compatibilidades que exige el paso a
profundidad arbitraria.
